\section*{Capítulo \ref{ch:b3:microbiomes} --- Microbiomas e simbioses}

\begin{solution}{exo:b3:microbiomes:1}
\omterm{def:b3:microbiomes:symbiosis}{Mutualismo}: os dois ganham --- leguminosa e rizóbio, coral e alga.
\omterm{def:b3:microbiomes:symbiosis}{Comensalismo}: um ganha e o outro não é afetado --- a maioria das
bactérias do intestino, que se alimentam do que o hospedeiro não
digere. Parasitismo: um ganha ao custo do outro --- a \emph{Wolbachia}
matando embriões machos. O mesmo par pode mudar: um comensal do
intestino que cruza a parede vira patógeno, um rizóbio que para de
fixar vira parasito do nódulo, e uma alga sob estresse térmico faz mal
ao coral que a abrigava.
\end{solution}

\begin{solution}{exo:b3:microbiomes:2}
$H = -(0.5\ln 0.5 + 0.3\ln 0.3 + 0.1\ln 0.1 + 2\times 0.05\ln 0.05) =
0.347 + 0.361 + 0.230 + 0.300 = 1.24$; $e^{H} = 3.4$ espécies
efetivas. $D = 0.25 + 0.09 + 0.01 + 0.0025 + 0.0025 = 0.355$; $1/D =
2.8$.
\end{solution}

\begin{solution}{exo:b3:microbiomes:3}
O levantamento do 16S relata quem está lá, até mais ou menos o gênero, e
suas contagens relativas de leituras; nada diz sobre os genes que eles
carregam, deixa passar fungos e \omterm{def:b3:virology:virus}{vírus} (que exigem outros \omterm{def:b3:genetic-engineering:pcr}{iniciadores}) e
é enviesado pelo número de cópias e pelos \omterm{def:b3:genetic-engineering:pcr}{iniciadores}. A metagenômica
aleatória relata genes, funções e, com profundidade bastante, genomas
e linhagens inteiras; custa mais por amostra, é soterrada pelo DNA do
hospedeiro em amostras de tecido, e ainda assim não consegue dizer que
genes estão expressos nem que células estão vivas.
\end{solution}

\begin{solution}{exo:b3:microbiomes:4}
Fermentar a fibra em \omterm{def:b3:microbiomes:gut}{ácidos graxos de cadeia curta} que alimentam o
cólon e o hospedeiro; sintetizar vitamina~K e vitaminas do complexo B;
transformar ácidos biliares e fármacos; a \omterm{def:b3:microbiomes:gut}{resistência à colonização}
contra patógenos; educar o sistema imune. Perturbação: a colite por
\emph{Clostridioides difficile} depois de \omterm{def:b3:bacteriology:antibiotics}{antibióticos} (também
implicadas: a doença inflamatória intestinal e a obesidade).
\end{solution}

\begin{solution}{exo:b3:microbiomes:5}
$\binom{20}{5} = \num{15504}$. Espécie com $10$: $1 - \binom{10}{5}/
\num{15504} = 1 - 252/\num{15504} = 0.984$; com $6$: $1 - 2002/\num{15504}
= 0.871$; com $3$: $1 - 6188/\num{15504} = 0.601$; com $1$: $1 -
\num{11628}/\num{15504} = 0.25$. $E[S_{5}] = 2.7$ espécies. Cobertura:
um singleton, $1 - 1/20 = 0.95$.
\end{solution}

\begin{solution}{exo:b3:microbiomes:6}
$400\times 10^{11} = 4\times 10^{13}$ bactérias, cerca de \qty{40}{g}.
Genes: $1000$ espécies $\times$ \num{4000}, descontada a sobreposição
--- da ordem de $3\times 10^{6}$ genes distintos, mais de cem vezes os
\num{20000} humanos.
\end{solution}

\begin{solution}{exo:b3:microbiomes:7}
A \omterm{def:b3:microbiomes:nodule}{nitrogenase} é destruída pelo oxigênio, mas os bacteroides precisam
respirar para fazer dezesseis ATP por N$_{2}$. O córtex do nódulo forma
uma barreira de difusão que limita a entrada de oxigênio, e a
\omterm{def:b3:microbiomes:nodule}{leghemoglobina}, em concentração milimolar, liga o oxigênio que entra,
de modo que a concentração livre fica perto de \qty{10}{nM} enquanto o
fluxo aos bacteroides é alto. O rosa é a oxileghemoglobina; um nódulo
verde degradou sua \omterm{def:b3:microbiomes:nodule}{leghemoglobina} --- ele é senescente e já não fixa.
\end{solution}

\begin{solution}{exo:b3:microbiomes:8}
Os dois falham: nem nódulos nem \omterm{def:b3:microbiomes:mycorrhiza}{micorriza arbuscular}, porque o pico de
cálcio é o sinal compartilhado da \omterm{def:b3:microbiomes:mycorrhiza}{via comum da simbiose}. A via,
portanto, é anterior à nodulação --- ela foi construída para a \omterm{def:b3:microbiomes:symbiosis}{simbiose}
fúngica quatrocentos milhões de anos atrás e recrutada pelas
leguminosas para os \omterm{def:b3:microbiomes:nodule}{rizóbios} mais tarde.
\end{solution}

\begin{solution}{exo:b3:microbiomes:9}
Kiers expôs alguns nódulos de soja a uma atmosfera de argônio e
oxigênio, na qual os \omterm{def:b3:microbiomes:nodule}{rizóbios} não conseguiam fixar nitrogênio, enquanto
os demais nódulos da planta fixavam normalmente; a planta cortou o
suprimento de oxigênio aos nódulos não fixadores, e os \omterm{def:b3:microbiomes:nodule}{rizóbios} deles
se reproduziram cerca da metade. Sem essa discriminação, os \omterm{def:b3:microbiomes:nodule}{rizóbios}
que pulassem a custosa fixação se reproduziriam mais depressa que os
fixadores, espalhar-se-iam pela população, e as plantas acabariam
abrigando bactérias que nada dariam: o \omterm{def:b3:microbiomes:symbiosis}{mutualismo} desmoronaria.
\end{solution}

\begin{solution}{exo:b3:microbiomes:10}
$\binom{N-N_{i}}{n}/\binom{N}{n} = \prod_{j=0}^{n-1}(N - N_{i} - j)/(N -
j)$, com cada fator abaixo de $1$; passar de $n$ a $n + 1$ multiplica
por mais um fator abaixo de $1$, de modo que a probabilidade de não ver
cai e a de ver sobe com $n$; em $n = N$ o numerador
$\binom{N -
N_{i}}{N}$ é zero e toda espécie está presente, dando $S$.
Para $N_{i},
n \ll N$ cada fator é de cerca de $1 - N_{i}/N$, e o
produto é de cerca de
$(1 -
N_{i}/N)^{n} \approx e^{-nN_{i}/N} \approx (1 - n/N)^{N_{i}}$ ---
a aproximação de amostragem com reposição.
\end{solution}

\begin{solution}{exo:b3:microbiomes:11}
Um macho é um beco sem saída para um simbionte que só viaja em ovos, de
modo que qualquer característica que transforme machos em fêmeas, que
os mate em favor de suas irmãs ou que prejudique fêmeas não infectadas
eleva a transmissão do simbionte. Incompatibilidade citoplasmática: o
esperma de machos infectados é modificado de modo que o zigoto morre a
menos que o ovo carregue a mesma \emph{Wolbachia}, que o resgata. As
fêmeas não infectadas, portanto, perdem a prole que concebem com machos
infectados, ao passo que as infectadas não perdem nenhuma; a vantagem
cresce com a frequência da infecção, de modo que, acima de um limiar, a
infecção se espalha até a fixação mesmo que as fêmeas infectadas ponham
um pouco menos ovos --- o custo é superado pela fração das ninhadas das
fêmeas não infectadas que é destruída.
\end{solution}

\begin{solution}{exo:b3:microbiomes:12}
O endossimbionte passa por um gargalo de algumas poucas células a cada
geração do hospedeiro e nunca recombina com forasteiros: a deriva fixa
mutações levemente deletérias, deleções inclusive, e não há fonte de
genes de reposição (a catraca de Muller). Os genes cujos produtos o
hospedeiro fornece, ou que servem a uma vida livre --- envoltório,
motilidade, regulação, a maior parte do reparo ---, não fazem falta e
vão primeiro; a maquinaria de tradução e os genes que fazem aquilo de
que o hospedeiro precisa (os aminoácidos essenciais na
\emph{Buchnera}) são mantidos por mais tempo. Os parentes de vida livre
têm populações enormes, recombinação e entrada de genes, e a seleção
mantém seus genomas. Reverter exigiria genes novos vindos de fora ---
impossível no isolamento ---, de modo que os hospedeiros, em vez disso,
substituem um simbionte erodido por um novo, como várias linhagens de
insetos já fizeram.
\end{solution}

\begin{solution}{pb:b3:microbiomes:1}
\textbf{1.} $400\times 10^{11} = 4\times 10^{13}$ bactérias,
\qty{40}{g}.
\textbf{2.} Um genoma de \num{4000} genes tem cerca de \qty{4}{Mb},
\qty{4.4}{fg}: $4\times 10^{13}\times 4.4\times 10^{-15} = \qty{0.18}{g}$ de DNA bacteriano, contra $3\times 10^{13}\times 6.4\times 10^{-12} = \qty{0.19}{g}$ de DNA humano --- praticamente
iguais.
\textbf{3.} $1000\times 0.7\times 4000 + 0.3\times 4000 \approx 2.8\times
10^{6}$ genes distintos, umas $140$ vezes os do genoma humano.
\textbf{4.} $30\times 10 = \qty{300}{mmol}$; $\times 0.9 = \qty{270}{kJ}$;
\qty{3}{\%} da dieta.
\textbf{5.} Não --- \qty{3}{\%} não explicam \qty{30}{\%}. Os
camundongos livres de germes também absorvem com menos eficiência, têm
hormônios intestinais, armazenamento de gordura e termogênese alterados
e apetite diferente: o \omterm{def:b3:microbiomes:symbiosis}{microbioma} age sobre a fisiologia do hospedeiro,
e não apenas como fornecedor de combustível.
\textbf{6.} Divisão equilibrada pela eliminação: uma população estável,
com cada célula reposta diariamente. Depois de uma morte de
\qty{99.9}{\%}, restam $4\times 10^{10}$ e dez duplicações restauram os
números em cerca de dez dias --- mas os sobreviventes e os de
crescimento mais rápido dominam, a composição fica alterada (uma
\omterm{def:b3:microbiomes:gut}{disbiose}) e, na brecha, um invasor como o \emph{C.~difficile} pode se
estabelecer.
\textbf{7.} As espécies menores dividem $0.4$ entre $177$: $p = 0.00226$
cada. $H
= -(0.3\ln 0.3 + 0.2\ln 0.2 + 0.1\ln 0.1 + 0.4\ln 0.00226) = 0.361 +
0.322 + 0.230 + 2.437 = 3.35$; $e^{H} \approx 28$.
\textbf{8.} $D = 0.09 + 0.04 + 0.01 + 177\times (0.00226)^{2} = 0.141$;
$1/D = 7.1$. O \omterm{def:b3:microbiomes:diversity}{índice de Simpson} eleva as abundâncias ao quadrado, de
modo que as três espécies dominantes o decidem e as raras quase não
contam; o de Shannon dá mais peso às espécies raras.
\textbf{9.} $1 - 60/2000 = 0.97$: cerca de $30$ leituras por mil
pertencem a espécies ainda não vistas.
\textbf{10.} A riqueza sobe com o tamanho da amostra enquanto espécies
raras continuam aparecendo; \num{10000} leituras alcançam espécies que
\num{2000} não alcançam. Rarefaça a amostra grande a \num{2000}
leituras pelo teorema (ou por subamostragem) --- ela deveria dar cerca
de $180$, se for a mesma comunidade --- e compare em profundidade
igual.
\textbf{11.} $N_{i} = 1$: $1 - \binom{1999}{1000}/\binom{2000}{1000} =
1 - (2000 - 1000)/2000 = 0.5$. $N_{i} = 5$: cerca de
$1 - 0.5^{5} = 0.97$.
\textbf{12.} $e^{2.0} = 7.4$ espécies efetivas; $D \ge 0.5^{2} = 0.25$,
cerca de $0.27$: uma comunidade com uma espécie dominante e algumas
dezenas de menores --- o intestino colapsado e de baixa diversidade em
que o \emph{C.~difficile} se instala.
\textbf{13.} $150\times 8 = \qty{1200}{kg}$ de açúcar por hectare,
\qty{10}{\%} das \qty{12}{t} fotossintetizadas.
\textbf{14.} $150\,000/28 = 5360$ mol de N$_{2}$; $16\times 5360 =
\num{85700}$ mol de ATP.
\textbf{15.} $\num{85700}/30 = 2860$ mol de glicose, \qty{514}{kg} ---
cerca de \qty{43}{\%} do açúcar pago. O resto constrói e mantém os
nódulos e os bacteroides, fornece os oito elétrons por N$_{2}$ como
redutor, e assimila a amônia.
\textbf{16.} $1200\times 16 = \qty{19200}{MJ}$ para \qty{150}{kg} de N:
\qty{128}{MJ/kg}, quase quatro vezes os \qty{35}{MJ/kg} da fábrica ---
mas pagos em luz do sol, no campo, sem combustível.
\textbf{17.} Fixadores $0.9\times 1 = 0.9$; trapaceiros $0.1\times 0.5 =
0.05$; os trapaceiros ficam em $0.05/0.95 = 5.3\,\%$ --- reduzidos à
metade em uma safra.
\textbf{18.} Trapaceiros $0.1\times 1.2^{10} = 0.62$ contra fixadores
$0.9$: \qty{41}{\%} depois de dez safras, rumo ao domínio. As sanções
voltam a seleção contra o trapaceiro; sem elas o \omterm{def:b3:microbiomes:symbiosis}{mutualismo} se erode.
\textbf{19.} $10^{6}\times \qty{20}{pg} = \qty{20}{\micro g}$ de carbono
por cm$^{2}$ por dia; \qty{18}{\micro g} para o coral.
\textbf{20.} Respiração \qty{15}{\micro g}: excedente de
\qty{3}{\micro g} por cm$^{2}$ por dia para crescimento, para a matriz
orgânica da calcificação, para muco e para gametas.
\textbf{21.} Com \qty{5}{\%} das algas, \qty{0.9}{\micro g} por dia
contra \qty{15}{\micro g} respirados: um déficit de \qty{14}{\micro g}
por dia; \qty{300}{\micro g} de reservas duram cerca de três semanas.
\textbf{22.} Quatro semanas a $+2$ é pior: o dano aos fotossistemas das
algas sobe abruptamente, e não linearmente, com a temperatura. A
métrica ainda assim funciona como aviso porque o branqueamento integra
o estresse cumulativo e os dois perfis diferem menos que a incerteza do
limiar; é um índice empírico, calibrado sobre branqueamentos
observados.
\textbf{23.} Um coral estressado pode deslocar-se rumo a espécies de
algas mais tolerantes ao calor já presentes em baixa abundância, ou
adquiri-las da água depois do branqueamento; as algas tolerantes dão
menos carbono, de modo que o coral cresce mais devagar, e a tolerância
ganha é de um ou dois graus --- menos que o aquecimento projetado ---,
enquanto o ritmo do aquecimento supera a adaptação dos próprios corais.
\textbf{24.} $10^{10}$ cm$^{2}$ $\times$ \qty{20}{\micro g} $=
\qty{200}{kg}$ de carbono por dia, cerca de \qty{73}{t} por ano.
\textbf{25.} Cerca de \qty{3}{\%} da dieta vindos da fermentação do
\omterm{def:b3:microbiomes:symbiosis}{microbioma}; \qty{97}{\%} de cobertura da amostra de sequenciamento;
\qty{10}{\%} da fotossíntese pagos pelo nitrogênio da lavoura de
feijão.
\end{solution}
